% ---
% title: An auditor's sample
% tags: finance, hypergeometric, expectation
% ---

\begin{problem}
	\begin{context}
		An auditor checks whether a company's accounts are correct. There are
		too many invoices to check every one, so the auditor picks a random
		sample and checks only those. An invoice is not picked twice: the
		sampling is \emph{without replacement}.
	\end{context}
	Of 20 invoices of a company, 4 contain an error. An auditor picks 5 of the
	20 invoices at random, without replacement. Let $X$ be the number of
	erroneous invoices in the sample.
	(a) What is the distribution of $X$, and what are its parameters?
	(b) Find $P(X=0)$, the probability that the auditor finds no error.
	(c) Find $P(X\ge1)$.
	(d) Find $P(X=2)$.
	(e) Find $\E[X]$.
\end{problem}

\begin{slideproblem}
	Of 20 invoices, 4 contain an error. An auditor picks 5 of them at random,
	without replacement; $X$ = number of erroneous invoices in the sample.
	(a) Distribution of $X$ and its parameters? \ (b) $P(X=0)$ (no error
	found) \ (c) $P(X\ge1)$ \ (d) $P(X=2)$ \ (e) $\E[X]$
\end{slideproblem}

\begin{slides}
	\small
	\exstatement
	\pause
	\begin{columns}[T]
		\begin{column}{.62\textwidth}
			\sol (a) $n=5$ of $N=20$ \alert{without replacement}, $m=4$
			erroneous: \alert{hypergeometric}, not binomial (draws not
			independent):
			$P(X=i)=\binom{4}{i}\binom{16}{5-i}\big/\binom{20}{5}$, with
			$\binom{20}{5}=15504$.\\[1mm]
			\pause
			(b) $P(X=0)=\binom{16}{5}\big/\binom{20}{5}=4368/15504\approx0.2817$.\\[1mm]
			\pause
			(c) Complement: $P(X\ge1)=1-P(X=0)\approx0.7183$.\\[1mm]
			\pause
			(d) $P(X=2)=\binom42\binom{16}{3}\big/\binom{20}{5}
				=6\cdot560/15504\approx0.2167$.\\[1mm]
			\pause
			(e) $\E[X]=nm/N=5\cdot4/20=1$.
		\end{column}
		\begin{column}{.36\textwidth}
			\centering
			\begin{tikzpicture}[x=.62cm,y=4.6cm]
			\draw[gray!40,very thin] (-0.6,0.1)--(4.6,0.1);
			\node[left,font=\tiny] at (-0.6,0.1) {$0.1$};
			\draw[gray!40,very thin] (-0.6,0.2)--(4.6,0.2);
			\node[left,font=\tiny] at (-0.6,0.2) {$0.2$};
			\draw[gray!40,very thin] (-0.6,0.3)--(4.6,0.3);
			\node[left,font=\tiny] at (-0.6,0.3) {$0.3$};
			\draw[gray!40,very thin] (-0.6,0.4)--(4.6,0.4);
			\node[left,font=\tiny] at (-0.6,0.4) {$0.4$};
			\draw[->] (-0.6,0)--(4.8999999999999995,0) node[right,font=\scriptsize] {$i$};
			\draw[->] (-0.6,0)--(-0.6,0.58) node[above,font=\scriptsize] {$P(X=i)$};
			\filldraw[fill=gold!55,draw=gold] (0-.38,0) rectangle (0+.38,0.2817);
			\node[below,font=\tiny] at (0,0) {$0$};
			\node[above,font=\tiny] at (0,0.2817) {(b)};
			\filldraw[fill=teal!22,draw=teal] (1-.38,0) rectangle (1+.38,0.4696);
			\node[below,font=\tiny] at (1,0) {$1$};
			\filldraw[fill=gold!55,draw=gold] (2-.38,0) rectangle (2+.38,0.2167);
			\node[below,font=\tiny] at (2,0) {$2$};
			\node[above,font=\tiny] at (2,0.2167) {(d)};
			\filldraw[fill=teal!22,draw=teal] (3-.38,0) rectangle (3+.38,0.031);
			\node[below,font=\tiny] at (3,0) {$3$};
			\filldraw[fill=teal!22,draw=teal] (4-.38,0) rectangle (4+.38,0.001);
			\node[below,font=\tiny] at (4,0) {$4$};
			\draw[navy,thick,dashed] (1,0)--(1,.54);
			\node[right,font=\tiny,text=navy] at (1,.54) {$\E[X]=1$};
			\end{tikzpicture}
		\end{column}
	\end{columns}
\end{slides}

\begin{solution}
	\idea The auditor takes a sample from a finite group of invoices without
	replacement, and counts the ``special'' (erroneous) ones in the sample. This
	is exactly the situation of the hypergeometric distribution.

	\step{(a)} We choose $n=5$ invoices from $N=20$, and $m=4$ of the 20 are
	erroneous. Every set of 5 invoices is equally likely, and no invoice can be
	chosen twice. So $X$, the number of erroneous invoices in the sample, is
	hypergeometric with parameters $N=20$, $m=4$, $n=5$:
	\[ P(X=i)=\frac{\dbinom{m}{i}\dbinom{N-m}{n-i}}{\dbinom{N}{n}}
		=\frac{\dbinom{4}{i}\dbinom{16}{5-i}}{\dbinom{20}{5}},
		\qquad i=0,1,2,3,4. \]
	In words: the denominator counts all possible samples; the numerator counts
	the samples with $i$ erroneous invoices (choose $i$ of the 4 erroneous ones)
	and $5-i$ correct invoices (choose them from the 16 correct ones). The
	number of all samples is
	\[ \binom{20}{5}=\frac{20\cdot19\cdot18\cdot17\cdot16}{5\cdot4\cdot3\cdot2\cdot1}=15504. \]

	\step{(b)} ``No error found'' means $X=0$: all 5 invoices come from the 16
	correct ones. There is $\binom40=1$ way to choose no erroneous invoice, and
	\[ \binom{16}{5}=\frac{16\cdot15\cdot14\cdot13\cdot12}{5\cdot4\cdot3\cdot2\cdot1}=4368 \]
	ways to choose the 5 correct ones. Hence
	\[ P(X=0)=\frac{1\cdot4368}{15504}=\frac{91}{323}\approx0.2817. \]
	So in more than one case out of four, the auditor finds nothing, even though
	a fifth of the invoices are wrong.

	\step{(c)} ``At least one'' would mean adding $P(X=1)$, $P(X=2)$, $P(X=3)$
	and $P(X=4)$: four terms. The complement of $X\ge1$ is $X=0$, which we
	already have, so it is shorter to use the complement rule:
	\[ P(X\ge1)=1-P(X=0)=1-\frac{91}{323}=\frac{232}{323}\approx0.7183. \]

	\step{(d)} For $X=2$ we choose 2 of the 4 erroneous invoices and 3 of the
	16 correct ones. Each choice of two erroneous invoices can be combined
	with every choice of three correct invoices, so the basic counting
	principle gives the product of the two counts:
	\[ \binom42=6,\qquad \binom{16}{3}=\frac{16\cdot15\cdot14}{3\cdot2\cdot1}=560, \]
	\[ P(X=2)=\frac{6\cdot560}{15504}=\frac{3360}{15504}=\frac{70}{323}\approx0.2167. \]

	\step{(e)} For a hypergeometric random variable $\E[X]=nm/N$:
	\[ \E[X]=\frac{5\cdot4}{20}=1. \]
	This is natural: a fifth of the invoices are erroneous ($m/N=4/20$), and
	the sample contains 5 invoices, so on average a fifth of 5, that is one
	erroneous invoice, is in the sample.
	\begin{answer}
		(a) hypergeometric, $N=20$, $m=4$, $n=5$ \quad
		(b) $P(X=0)=\frac{4368}{15504}=\frac{91}{323}\approx0.2817$ \quad
		(c) $P(X\ge1)\approx0.7183$ \quad
		(d) $P(X=2)=\frac{3360}{15504}=\frac{70}{323}\approx0.2167$ \quad
		(e) $\E[X]=1$
	\end{answer}
	\begin{remark}
		\textbf{Common mistake:} using the binomial distribution with $n=5$ and
		$p=4/20=0.2$. The binomial needs independent trials with the same success
		probability. Here the draws are not independent: given that the first
		invoice is erroneous, the second is erroneous with probability $3/19$;
		given that the first is correct, the probability is $4/19$.
		Unconditionally it is $4/20$. This is
		what ``without replacement'' changes. The binomial would give
		$P(X=0)=0.8^5\approx0.3277$ instead of $0.2817$. The expected value is
		the same in both models ($np=5\cdot0.2=1$), but the probabilities differ.
	\end{remark}
\end{solution}
